\begin{theorem}[奇异环椭圆平均的壁穿越公式：有理测试函数的离散跳跃律]\label{thm:xi-singular-ring-ellipse-wallcrossing-rational}
沿用命题 \ref{prop:xi-joukowsky-ellipse-average-constant-term} 的记号，并令 \(r:=\sqrt{L}\)。
对任意有理函数 \(f\in\CC(w)\) 且满足 \(f\) 在 \(E_L\) 上无极点，定义奇异环椭圆平均
\[
\mathcal{L}_L(f):=\int_{E_L} f(S)\,\dd\nu_L(S)
=\frac{1}{2\pi}\int_{0}^{2\pi} f\!\left(r e^{\mathrm{i}\theta}+r^{-1}e^{-\mathrm{i}\theta}\right)\dd\theta.
\]
令
\[
g(z):=f(z+z^{-1})\in\CC(z),
\qquad
h(z):=\frac{g(z)}{z}.
\]
若 \(r\) 不等于 \(h\) 的任一非零极点模长，则有留数表示
\[
\boxed{\;
\mathcal{L}_L(f)=\Res_{z=0}h(z)\;+\;\sum_{\substack{a\neq 0\\ a\ \text{为 }h\text{ 的极点}\\ 0<|a|<r}}\Res_{z=a}h(z).\;}
\]
并且将 \(t:=\log r=\tfrac12\log L\) 视为参数，则 \(\mathcal{L}_L(f)\) 在去除离散集合
\(\{\log|a|:a\neq 0,\ a\ \text{为 }h\text{ 的极点}\}\)
后对 \(t\) 局部常值；当 \(r\) 穿越某一极点 \(a\neq 0\) 的模长时，其跳跃量严格等于 \(\Res_{z=a}h(z)\)。
\end{theorem}
\begin{proof}
以 \(\zeta=e^{\mathrm{i}\theta}\) 代换得 \(\dd\theta=\dd\zeta/(\mathrm{i}\zeta)\)，从而
\[
\mathcal{L}_L(f)=\frac{1}{2\pi\mathrm{i}}\oint_{|\zeta|=1} f(r\zeta+r^{-1}\zeta^{-1})\,\frac{\dd\zeta}{\zeta}.
\]
作代换 \(z=r\zeta\)（于是 \(\dd\zeta/\zeta=\dd z/z\)）得到
\[
\mathcal{L}_L(f)=\frac{1}{2\pi\mathrm{i}}\oint_{|z|=r} f(z+z^{-1})\,\frac{\dd z}{z}
=\frac{1}{2\pi\mathrm{i}}\oint_{|z|=r} h(z)\,\dd z.
\]
当 \(r\) 的微扰不穿越任何极点时，可在不越过极点的条件下连续变形围道，故 \(\mathcal{L}_L(f)\) 关于 \(t\) 局部常值。
由留数定理，上式积分等于圆盘 \(|z|<r\) 内全部留数之和；其中 \(z=0\) 的留数恒计入，而其余留数仅在相应极点落入 \(|z|<r\) 时出现，从而得到所示表示式，并给出穿越模长时的跳跃律。证毕。
\end{proof}

\begin{proposition}[有限原子跳跃谱的椭圆平均实现：给定跳跃数据构造有理测试函数]\label{prop:xi-ellipse-wallcrossing-realize-finite-jump-data}
\leanverified{paper\_xi\_ellipse\_wallcrossing\_realize\_finite\_jump\_data}
取有限集合 \(\{a_1,\dots,a_N\}\subset(1,\infty)\) 与权重 \(c_1,\dots,c_N\in\CC\)，并定义
\[
g(z):=\sum_{j=1}^N c_j\left(\frac{z}{z-a_j}+\frac{1}{1-a_j z}\right)\in\CC(z).
\]
则 \(g(z)=g(1/z)\)，因而存在唯一 \(f\in\CC(w)\) 满足 \(g(z)=f(z+z^{-1})\)。
并且对任意 \(r>1\) 且 \(r\notin\{a_1,\dots,a_N\}\) 有
\[
\boxed{\;
\mathcal{L}_{r^2}(f)=\sum_{j:\ a_j<r} c_j\; }.
\]
换言之，\(\mathcal{L}_{r^2}(f)\) 作为 \(r\) 的函数为分段常值并仅在 \(r=a_j\) 处发生跳跃，且跳跃幅度恰为 \(c_j\)。
\end{proposition}
\begin{proof}
逐项通分可核对对每个 \(a>1\) 都有
\(
\frac{z}{z-a}+\frac{1}{1-a z}
\)
在变换 \(z\mapsto 1/z\) 下不变，从而 \(g(1/z)=g(z)\)。
由于 \(\CC(z)^{\langle z\mapsto 1/z\rangle}=\CC(z+z^{-1})\)，存在唯一 \(f\in\CC(w)\) 使 \(g(z)=f(z+z^{-1})\)。
\par
令 \(h(z):=g(z)/z\)。
按定义，\(h\) 在 \(z=a_j\) 处的留数为 \(c_j\)，并在 \(z=1/a_j\) 处的留数为 \(-c_j\)。
当 \(r>1\) 时总有 \(|1/a_j|<1<r\)，故所有 \(z=1/a_j\) 的留数恒计入圆盘 \(|z|<r\)，而 \(z=a_j\) 的留数当且仅当 \(a_j<r\) 时计入。
另一方面 \(\Res_{z=0}h(z)=\lim_{z\to 0}g(z)=\sum_{j=1}^N c_j\)。
将这些量代入定理 \ref{thm:xi-singular-ring-ellipse-wallcrossing-rational} 的留数公式即得
\(
\mathcal{L}_{r^2}(f)=\sum_{j:\ a_j<r}c_j
\)。
证毕。
\end{proof}
\begin{corollary}[奇异环椭圆平均的单位跳跃普适性：任意有限 $0/1$ 阶梯的留数实现]\label{cor:xi-ellipse-wallcrossing-unit-staircase-universality}
\leanverified{paper\_xi\_ellipse\_wallcrossing\_unit\_staircase\_universality}
沿用定理 \ref{thm:xi-singular-ring-ellipse-wallcrossing-rational} 的记号，并令 \(r=\sqrt{L}\)、\(t=\log r\)。
给定任意有限集合 \(S=\{\delta_1,\dots,\delta_M\}\subset\RR_{>0}\)，存在显式有理函数 \(f_S\in\CC(w)\)（其极点均不落在 \(E_L\) 上，除去离散的 \(L\) 集外）使得对所有 \(t\in\RR\) 都有分段常值恒等式
\[
\mathcal{L}_{r=e^{t}}(f_S)+M
=
\#\{\delta\in S:\ \delta<t\}.
\]
从而 \(\mathcal{L}_{r=e^t}(f_S)\) 在 \(t\)-轴上的每一次跳跃都为严格的单位跳跃，且跳点集合恰为 \(S\)。
\end{corollary}
\begin{proof}
对每个 \(\delta\in S\) 令 \(a_\delta:=e^{\delta}>1\)，并定义
\[
f_\delta(w):=\frac{a_\delta-a_\delta^{-1}}{w-(a_\delta+a_\delta^{-1})},
\qquad
f_S:=\sum_{\delta\in S} f_\delta.
\]
对单个 \(f_\delta\)，令 \(g_\delta(z)=f_\delta(z+z^{-1})\)、\(h_\delta(z)=g_\delta(z)/z\)。
直接化简得
\[
h_\delta(z)=\frac{a_\delta-a_\delta^{-1}}{(z-a_\delta)(z-a_\delta^{-1})},
\]
其极点为 \(a_\delta\) 与 \(a_\delta^{-1}\)，且
\[
\Res_{z=a_\delta}h_\delta=1,\qquad \Res_{z=a_\delta^{-1}}h_\delta=-1.
\]
当 \(r>1\) 时，\(|a_\delta^{-1}|<1<r\) 恒成立，因此极点 \(a_\delta^{-1}\) 的留数 \(-1\) 始终计入定理 \ref{thm:xi-singular-ring-ellipse-wallcrossing-rational} 的求和；
而极点 \(a_\delta\) 被计入当且仅当 \(r>|a_\delta|=e^\delta\)，并额外贡献 \(+1\)。
故对 \(r>1\) 且 \(r\neq e^\delta\) 有
\[
\mathcal{L}_L(f_\delta)=-1+\mathbf 1_{\{r>e^\delta\}}.
\]
对 \(\delta\in S\) 求和得到
\[
\mathcal{L}_L(f_S)=-M+\#\{\delta\in S:\ r>e^\delta\}.
\]
令 \(r=e^t\) 即得结论。证毕。
\end{proof}


\begin{corollary}[单极核的完全显式分段常值：\texorpdfstring{$(w-a)^{-1}$}{(w-a)^-1} 的奇异环响应]\label{cor:xi-singular-ring-ellipse-simple-pole-response}
取 \(f(w)=\frac{1}{w-a}\)，并假设 \(a\notin E_L\)。
令
\[
z_\pm(a):=\frac{a\pm\sqrt{a^2-4}}{2},
\qquad
\text{其中 }\sqrt{\cdot}\text{取任意固定分支且 }z_+(a)z_-(a)=1.
\]
则对任意 \(r=\sqrt{L}>0\) 且 \(r\neq |z_\pm(a)|\) 有
\[
\boxed{\;
\mathcal{L}_L\!\left(\frac{1}{w-a}\right)=
\begin{cases}
0, & 0<r<\min\{|z_+(a)|,|z_-(a)|\},\\[4pt]
-\dfrac{1}{\sqrt{a^2-4}}, & \min\{|z_+(a)|,|z_-(a)|\}<r<\max\{|z_+(a)|,|z_-(a)|\},\\[8pt]
0, & r>\max\{|z_+(a)|,|z_-(a)|\}.
\end{cases}}
\]
因此该核的椭圆平均在 \(\log r\) 方向上至多发生一次跳跃，其幅度为 \(\pm 1/\sqrt{a^2-4}\)（符号由平方根分支决定）。
\end{corollary}
\begin{proof}
由恒等式
\[
\frac{1}{z+z^{-1}-a}\cdot\frac{1}{z}=\frac{1}{z^2-a z+1}
\]
知 \(h(z)=\frac{1}{z^2-a z+1}\) 的极点为 \(z_\pm(a)\)，且
\[
\Res_{z=z_\pm(a)}\frac{1}{z^2-a z+1}
=\frac{1}{2z_\pm(a)-a}
=\pm\frac{1}{\sqrt{a^2-4}}.
\]
由定理 \ref{thm:xi-singular-ring-ellipse-wallcrossing-rational}，\(\mathcal{L}_L(f)\) 等于 \(|z|<r\) 内极点留数之和。
随 \(r\) 增大，极点按模长进入圆盘；当两极点均落入 \(|z|<r\) 时两留数相加为零，故仅在“恰有一极点进入”区间出现非零值，从而得到分段公式。证毕。
\end{proof}

\begin{theorem}[Joukowsky 奇异环下的跳跃互换律：椭圆 Cauchy 跳跃的平方根分母]\label{thm:xi-joukowsky-ellipse-plemelj-jump}
令 \(L>1\)，并设 \(\sigma\) 为 \(\TT\) 上复测度且在角变量 \(\theta\) 下绝对连续：
\[
\dd\sigma(e^{\mathrm{i}\theta})=\rho(\theta)\frac{\dd\theta}{2\pi},
\qquad \rho\in L^p([0,2\pi]),\ p>1.
\]
令 \(\sigma_L:=(J_L)_\#\sigma\) 为其在 \(E_L\) 上的推前测度，并定义 \(E_L\) 上的 Cauchy 变换
\[
G_{\sigma_L}(w):=\int_{E_L}\frac{1}{w-\xi}\,\dd\sigma_L(\xi),\qquad w\in\CC\setminus E_L.
\]
则对几乎处处的 \(\theta\)，沿椭圆参数化 \(w(\theta):=J_L(e^{\mathrm{i}\theta})\) 的两侧非切向边界值存在，并满足跳跃恒等式
\[
\boxed{\;
G_{\sigma_L}^{\mathrm{ext}}\!\bigl(w(\theta)\bigr)-G_{\sigma_L}^{\mathrm{int}}\!\bigl(w(\theta)\bigr)
=-\frac{\rho(\theta)}{\sqrt{w(\theta)^2-4}}.\;}
\]
其中平方根分支取为椭圆上满足
\[
\sqrt{w(\theta)^2-4}=r e^{\mathrm{i}\theta}-r^{-1}e^{-\mathrm{i}\theta},
\qquad r=\sqrt{L},
\]
并与无穷远处的渐近 \(\sqrt{w^2-4}\sim w\) 连续相容者。
\end{theorem}
\begin{proof}
定理 \ref{thm:xi-joukowsky-cauchy-transform-two-branch} 给出对 \(w\notin E_L\) 的双分支表示
\[
G_{\sigma_L}(w)=\frac{1}{\sqrt{w^2-4}}
\Bigl(z_-(w)\,G_\sigma(z_-(w)) - z_+(w)\,G_\sigma(z_+(w))\Bigr),
\]
其中 \(z_\pm(w)=\frac{w\pm\sqrt{w^2-4}}{2\sqrt{L}}\)，且
\(
G_\sigma(z)=\int_{\TT}\frac{1}{\zeta-z}\,\dd\sigma(\zeta)
\)。
在上述平方根分支下，当 \(w\to w(\theta)\) 时有
\[
z_+(w(\theta))=e^{\mathrm{i}\theta},
\qquad
z_-(w(\theta))=\frac{e^{-\mathrm{i}\theta}}{L},
\]
从而 \(z_-(w(\theta))\) 远离单位圆（当 \(L\neq 1\) 时 \(|z_-(w(\theta))|=1/L\neq 1\)），故相应项不产生边界跳跃。

另一方面，把 \(\sigma\) 写为复形式 \( \dd\sigma(\zeta)=f(\zeta)\,\dd\zeta \)（\(\zeta=e^{\mathrm{i}\theta}\)）：
\[
f(\zeta)=\frac{\rho(\theta)}{2\pi\mathrm{i}\,\zeta}.
\]
由 \(p>1\) 的 Plemelj 定理，几乎处处 \(\theta\) 上成立
\[
G_\sigma^{\mathrm{ext}}(e^{\mathrm{i}\theta})-G_\sigma^{\mathrm{int}}(e^{\mathrm{i}\theta})
=2\pi\mathrm{i}\,f(e^{\mathrm{i}\theta})
=\frac{\rho(\theta)}{e^{\mathrm{i}\theta}}.
\]
将其代入双分支表示式的边界极限，并注意 \(z_+(w(\theta))=e^{\mathrm{i}\theta}\)，得到
\[
G_{\sigma_L}^{\mathrm{ext}}(w(\theta))-G_{\sigma_L}^{\mathrm{int}}(w(\theta))
=-\frac{1}{\sqrt{w(\theta)^2-4}}\cdot e^{\mathrm{i}\theta}\Bigl(G_\sigma^{\mathrm{ext}}(e^{\mathrm{i}\theta})-G_\sigma^{\mathrm{int}}(e^{\mathrm{i}\theta})\Bigr)
=-\frac{\rho(\theta)}{\sqrt{w(\theta)^2-4}}.
\]
证毕。
\end{proof}

\begin{theorem}[“二值跳跃”的统一表述：取向函子同时控制解析跳跃与离散 uplift]\label{thm:xi-binary-jump-orientation-unification}
令 \(\pi:\widetilde Y\to Y\) 为一族局部有限层覆叠，使得对每个 \(y\in Y\) 的纤维 \(S_y:=\pi^{-1}(y)\) 为有限集合；
允许的“重标号”由 \(\Aut(S_y)\cong \mathfrak S_{|S_y|}\) 给出。
对每个纤维定义其取向 \(\ZZ_2\)-torsor
\[
\mathrm{Or}(S_y):=\mathrm{Bij}([|S_y|],S_y)\big/\mathfrak A_{|S_y|}.
\]
则在本文所出现的两类典型二值跳跃机制中，二值对象都可自然表达为 \(\mathrm{Or}\) 的像：
\begin{enumerate}
\normalfont
  \item \textbf{离散 uplift：}对 boundary uplift 位移集合 \(\Delta_m^{\mathrm{bdry}}(w)\)（作为无序有限集），在不引入外加寄存器且要求置换自然性的约束下，唯一的非平凡二值跳跃对象可取为
  \(\mathrm{Or}(\Delta_m^{\mathrm{bdry}}(w))\)（推论 \ref{cor:bdry-orientation-parity-uplift}）。
  \item \textbf{解析跳跃：}在 Joukowsky 二重覆盖的两分支表示中（定理 \ref{thm:xi-joukowsky-cauchy-transform-two-branch}），对椭圆边界点 \(w(\theta)=J_L(e^{\mathrm{i}\theta})\) 存在两侧边界值，
  且 Cauchy 跳跃
  \(
  G_{\sigma_L}^{\mathrm{ext}}(w(\theta))-G_{\sigma_L}^{\mathrm{int}}(w(\theta))
  \)
  在交换两支（等价地在双覆盖纤维 \(\{z_+(w(\theta)),z_-(w(\theta))\}\) 上作用非平凡置换）下变号（定理 \ref{thm:xi-joukowsky-ellipse-plemelj-jump} 的平方根分支互换律）。
因而该跳跃量可视为取向 torsor \(\mathrm{Or}(\{z_+(w(\theta)),z_-(w(\theta))\})\) 的符号等变截面，从而不依赖对两支的标号选择。
\end{enumerate}
因此，在“允许重标号但禁止外加寄存器”的共同口径下，二值跳跃的可审计结构由取向函子统一刻画；并且在局部层数 \(\ge 3\) 的离散情形中，该二值对象在同构意义下由符号表示唯一确定（定理 \ref{thm:bdry-binary-jump-orientation-functor-uniqueness}）。
\leanverified{paper\_xi\_binary\_jump\_orientation\_unification}
\end{theorem}

\begin{proof}
离散部分由推论 \ref{cor:bdry-orientation-parity-uplift} 直接给出。
解析部分仅用对称性：二重覆盖的纤维为二元集合，其非平凡置换生成 \(\mathfrak S_2\)；
而跳跃互换律表明所述差量对该置换取负，故其变换律正是符号表示，从而可在取向 torsor 上无歧义地表达为符号等变截面。
最后一段的“局部层数 \(\ge 3\) 唯一性”由定理 \ref{thm:bdry-binary-jump-orientation-functor-uniqueness} 的分类结论给出。证毕。
\end{proof}

\begin{conjecture}[window-\(6\) 二层覆盖的函子化分支：由 sheet parity 到 \texorpdfstring{$\mathrm{Spin}(10)$}{Spin(10)} 的半自旋选择]\label{conj:xi-window6-sheet-parity-functor-spin10-branching}
在 boundary uplift 的口径下，假定 window \(m=6\) 满足如下二层刚性：对每个边界词 \(w\)，其微观 uplift 纤维恰含两个预像，并且在一组固定编号制度下两预像的编号差为常数
\(\delta=34(=F_9)\)；
由此得到规范的 \(\ZZ_2\)-分级（sheet parity）。
令 \(\mathcal B_6\) 为由该二层覆盖数据生成的群胚/覆叠范畴（对象为边界词，态射为与 uplift 相容的离散变换，并要求置换自然性且不引入外加寄存器）。
则存在自然的 \(\ZZ_2\)-分级函子
\[
\Xi:\ \mathcal B_6\longrightarrow \mathrm{Rep}\!\bigl(\mathrm{Spin}(10)\bigr),
\]
使得 \(\Xi\) 在 \(D_5\) 的外自同构类 \(\mathrm{Out}(D_5)\cong\ZZ_2\) 上取到非平凡像，从而在两支半自旋表示之间给出与 fold-aware restriction 相容的唯一分支选择。
并且当 \(m\neq 6\) 时（例如 \(m=7,8\) 的三层 uplift），上述意义下的稳定 \(\ZZ_2\)-分级不再存在，因此不存在与之相容的同构意义上的“手征选择律”。
\end{conjecture}

\begin{theorem}[布尔二值跳跃核的张量谱与互补反演]\label{thm:xi-boolean-binary-jump-kernels-tensor-spectrum}
固定整数 \(q\ge 1\)，并记 \(\Omega_q:=\{0,1\}^q\)。
对 \(u,v\in\Omega_q\) 记按位运算 \(u\vee v\)、\(u\wedge v\)，并记 \(\mathbf{1}=(1,\dots,1)\)、\(\mathbf{0}=(0,\dots,0)\)。
定义两类二值跳跃核（大小 \(2^q\times 2^q\)）：
\[
\mathsf{C}_q(u,v):=\mathbf{1}_{\{u\vee v=\mathbf{1}\}},
\qquad
\mathsf{D}_q(u,v):=\mathbf{1}_{\{u\wedge v=\mathbf{0}\}}.
\]
令
\[
A:=\begin{pmatrix}0&1\\[2pt]1&1\end{pmatrix},
\qquad
B:=\begin{pmatrix}1&1\\[2pt]1&0\end{pmatrix},
\qquad
\varphi:=\frac{1+\sqrt5}{2}.
\]
则：
\begin{enumerate}
  \item \textbf{张量分解：}在标准基下有 \(\mathsf{C}_q=A^{\otimes q}\)、\(\mathsf{D}_q=B^{\otimes q}\)。
  \item \textbf{黄金比例谱闭式：}\(\mathsf{C}_q\) 与 \(\mathsf{D}_q\) 具有同一特征值多重集：
  \[
  \lambda_k^{(q)}=\varphi^k(-\varphi^{-1})^{\,q-k}=(-1)^{q-k}\varphi^{\,2k-q}\qquad(k=0,1,\dots,q),
  \]
  且每个 \(\lambda_k^{(q)}\) 的重数为 \(\binom{q}{k}\)。
  特别地，谱半径满足 \(\rho(\mathsf{C}_q)=\rho(\mathsf{D}_q)=\varphi^q\)。
  \item \textbf{行列式与整模性：}
  \[
  \det(\mathsf{C}_q)=\det(\mathsf{D}_q)=(-1)^{q2^{q-1}}.
  \]
  因而当 \(q\ge 2\) 时二者为幺模矩阵，逆矩阵为整数矩阵。
  \item \textbf{逆核的显式支持：}在按坐标分解的乘积基下有
  \[
  \mathsf{C}_q^{-1}(u,v)=(-1)^{\#\{i:\ u_i=v_i=0\}}\mathbf{1}_{\{u\wedge v=\mathbf{0}\}},
  \qquad
  \mathsf{D}_q^{-1}(u,v)=(-1)^{|u\wedge v|}\mathbf{1}_{\{u\vee v=\mathbf{1}\}}.
  \]
  \item \textbf{互补共轭：}令 \(\kappa:\Omega_q\to\Omega_q\) 为按位补集 \(\kappa(u)=\mathbf{1}-u\)，并令 \(P_\kappa\) 为其置换矩阵。
  则 \(\mathsf{C}_q=P_\kappa\,\mathsf{D}_q\,P_\kappa^{\mathsf T}\)，从而两者谱完全一致。
\end{enumerate}
\end{theorem}
\leanverified{paper\_xi\_boolean\_binary\_jump\_kernels\_tensor\_spectrum}
\begin{proof}
对 \(q=1\) 直接核对得到 \(A_{u,v}=\mathbf{1}_{\{u\vee v=1\}}\)、\(B_{u,v}=\mathbf{1}_{\{u\wedge v=0\}}\)。
对一般 \(q\)，由逐坐标分解
\(
\mathbf{1}_{\{u\vee v=\mathbf{1}\}}=\prod_{i=1}^q \mathbf{1}_{\{u_i\vee v_i=1\}}
\)
与
\(
\mathbf{1}_{\{u\wedge v=\mathbf{0}\}}=\prod_{i=1}^q \mathbf{1}_{\{u_i\wedge v_i=0\}}
\)
可得张量分解。
\par
矩阵 \(A,B\) 具有同一特征多项式 \(\lambda^2-\lambda-1\)，故特征值为 \(\varphi\) 与 \(-\varphi^{-1}\)。
Kronecker 幂的谱为基矩阵谱的乘积多重集，取 \(\varphi\) 的次数为 \(k\) 的组合数为 \(\binom{q}{k}\)，从而得到谱闭式。
\par
行列式公式由 \(\det(M\otimes N)=\det(M)^{\dim N}\det(N)^{\dim M}\) 迭代得到，故 \(\det(M^{\otimes q})=\det(M)^{q2^{q-1}}\)；而 \(\det(A)=\det(B)=-1\)，即得结论。
当 \(q\ge 2\) 时行列式为 \(+1\)，从而逆为整数矩阵。
\par
最后，利用 \((M^{\otimes q})^{-1}=(M^{-1})^{\otimes q}\) 与
\(
A^{-1}=\begin{psmallmatrix}-1&1\\[2pt]1&0\end{psmallmatrix},
\;
B^{-1}=\begin{psmallmatrix}0&1\\[2pt]1&-1\end{psmallmatrix}
\)
逐坐标相乘即可得到逆核的支持条件与符号因子。
互补共轭由恒等式
\(
(1-u_i)\wedge(1-v_i)=0\iff u_i\vee v_i=1
\)
逐坐标推出。证毕。
\end{proof}

\begin{theorem}[特征 \(2\) 下布尔二值跳跃核的三阶周期性与 \texorpdfstring{$S_3$}{S3}-结构]\label{thm:xi-boolean-binary-jump-kernels-char2-s3}
沿用定理 \ref{thm:xi-boolean-binary-jump-kernels-tensor-spectrum} 的记号，并令
\[
\overline{\mathsf C}_q,\overline{\mathsf D}_q\in \mathrm{GL}(2^q,\FF_2)
\]
分别为 \(\mathsf C_q,\mathsf D_q\) 在 \(\FF_2\) 上的约化。
则对任意 \(q\ge 1\)：
\begin{enumerate}
\normalfont
  \item \textbf{三阶周期性：}\(\overline{\mathsf C}_q^3=I\) 且 \(\overline{\mathsf D}_q^3=I\)。
  \item \textbf{互为逆与交换：}\(\overline{\mathsf D}_q=\overline{\mathsf C}_q^{\,2}=\overline{\mathsf C}_q^{-1}\)，并且
  \(
  \overline{\mathsf C}_q\,\overline{\mathsf D}_q=I=\overline{\mathsf D}_q\,\overline{\mathsf C}_q
  \)。
  \item \textbf{互补共轭生成 \texorpdfstring{$S_3$}{S3}：}
  令 \(P_\kappa\) 为按位互补 \(\kappa(u)=\mathbf 1-u\) 的置换矩阵（定理 \ref{thm:xi-boolean-binary-jump-kernels-tensor-spectrum}）。
  则
  \[
  P_\kappa^2=I,\qquad
  P_\kappa\,\overline{\mathsf C}_q\,P_\kappa=\overline{\mathsf C}_q^{-1},
  \]
  从而
  \(
  \langle \overline{\mathsf C}_q,P_\kappa\rangle\simeq D_6\simeq S_3
  \)。
  \item \textbf{谱型与不可约分解：}
  设 \(\FF_4/\FF_2\) 为分裂扩张，并取本原三次单位根 \(\omega\in\FF_4^\times\)（\(\omega^2+\omega+1=0\)）。
  则 \(\overline{\mathsf C}_q\) 在 \(\FF_4\) 上可对角化，其特征值仅为 \(1,\omega,\omega^2\)，且重数满足
  \[
  m_0:=\dim_{\FF_2}\ker(\overline{\mathsf C}_q-I)=\frac{2^q+2(-1)^q}{3},
  \qquad
  m_1=m_2=\frac{2^q-(-1)^q}{3}.
  \]
  等价地，作为 \(\FF_2[C_3]\)-模有分解
  \[
  \FF_2^{\Omega_q}\ \cong\ \mathbf 1^{\oplus m_0}\ \oplus\ V^{\oplus m},
  \qquad
  m=\frac{2^q-(-1)^q}{3},
  \]
  其中 \(\mathbf 1\) 为平凡一维表示，而 \(V\) 为 \(\FF_2[C_3]\) 的唯一二维不可约表示（对应不可约多项式 \(t^2+t+1\)）。
\end{enumerate}
\end{theorem}
\begin{proof}
由定理 \ref{thm:xi-boolean-binary-jump-kernels-tensor-spectrum}，有 \(\mathsf C_q=A^{\otimes q}\)、\(\mathsf D_q=B^{\otimes q}\) 且 \(B=A^2\)。
将 \(A\) 约化至 \(\FF_2\) 得
\(
\overline A=\begin{psmallmatrix}0&1\\1&1\end{psmallmatrix}
\),
并直接计算得 \(\overline A^3=I\) 与 \(\overline A^2=\overline B\)。
因此
\[
\overline{\mathsf C}_q^3=(\overline A^{\otimes q})^3=(\overline A^3)^{\otimes q}=I,
\qquad
\overline{\mathsf D}_q=\overline B^{\otimes q}=(\overline A^2)^{\otimes q}=\overline{\mathsf C}_q^{\,2},
\]
从而 (1)(2) 成立。
\par
对 (3)，由定理 \ref{thm:xi-boolean-binary-jump-kernels-tensor-spectrum} 的互补共轭恒等式
\(
\mathsf C_q=P_\kappa \mathsf D_q P_\kappa^{\mathsf T}
\)
在 \(\FF_2\) 上约化得到
\(
P_\kappa\,\overline{\mathsf C}_q\,P_\kappa=\overline{\mathsf D}_q=\overline{\mathsf C}_q^{-1}
\)。
结合 \(P_\kappa^2=I\) 与 \(\overline{\mathsf C}_q^3=I\)，得到二面体关系 \(srs=r^{-1}\)（\(r=\overline{\mathsf C}_q\)，\(s=P_\kappa\)），故生成群同构于 \(D_6\simeq S_3\)。
\par
对 (4)，在 \(\FF_4\) 上 \(\overline A\) 的特征多项式为 \(t^2+t+1=(t-\omega)(t-\omega^2)\)，故 \(\overline A\) 相似于 \(\mathrm{diag}(\omega,\omega^2)\) 并可对角化。
于是 \(\overline{\mathsf C}_q=\overline A^{\otimes q}\) 的特征值为
\(
\omega^{\epsilon_1+\cdots+\epsilon_q}
\)
其中每个 \(\epsilon_i\in\{1,2\}\)。
令 \(m_r\) 表示特征值 \(\omega^r\)（\(r=0,1,2\)）的重数，则 \(m_r\) 等于满足
\(
\epsilon_1+\cdots+\epsilon_q\equiv r\ (\mathrm{mod}\ 3)
\)
的序列数。
令 \(\zeta:=e^{2\pi\mathrm i/3}\) 为复三次单位根。
单位根滤波给出
\[
m_r=\frac{1}{3}\sum_{t=0}^{2}\zeta^{-rt}\bigl(\zeta^t+\zeta^{2t}\bigr)^q
\;=\;
\frac{1}{3}\Bigl(2^q+(-1)^q(\zeta^{-r}+\zeta^{-2r})\Bigr),
\]
其中使用 \(\zeta+\zeta^2=-1\)。
取 \(r=0\) 得 \(m_0=(2^q+2(-1)^q)/3\)；
而 \(r=1,2\) 时 \(\omega^{-r}+\omega^{-2r}=\omega+\omega^2=-1\)，故 \(m_1=m_2=(2^q-(-1)^q)/3\)。
\par
最后，由 \(\FF_2\) 上多项式分解
\(
t^3-1=(t-1)(t^2+t+1)
\)
与 \(t^2+t+1\) 的不可约性可知：\(\overline{\mathsf C}_q\) 的 \(1\)-特征空间给出 \(m_0\) 个平凡分量，而其余部分在 \(\FF_2\) 上必分解为 \(m_1=m_2\) 个二维不可约分量，记为 \(V\)。
从而得到所示 \(\FF_2[C_3]\)-模分解。证毕。
\end{proof}
\leanverified{paper\_xi\_boolean\_binary\_jump\_kernels\_char2\_s3}

\leanverified{paper\_xi\_boolean\_binary\_jump\_kernels\_dihedral\_inversion}
\begin{proposition}[逆核的 Hamming--符号对角相似与二面体反演]\label{prop:xi-boolean-binary-jump-kernels-dihedral-inversion}
沿用定理 \ref{thm:xi-boolean-binary-jump-kernels-tensor-spectrum} 的记号。
对 \(u\in\Omega_q\) 记 \(|u|\) 为 Hamming 重，并定义对角符号矩阵
\[
H:=\mathrm{diag}\bigl((-1)^{|u|}\bigr)_{u\in\Omega_q}\in \mathrm{GL}(2^q,\ZZ).
\]
则有严格恒等式
\[
\mathsf D_q^{-1}=(-1)^q\,H\,\mathsf C_q\,H,
\qquad
\mathsf C_q^{-1}=(-1)^q\,H\,\mathsf D_q\,H.
\]
并且对互补置换矩阵 \(P_\kappa\) 有
\[
\mathsf C_q^{-1}=(-1)^q\,(H P_\kappa)\,\mathsf C_q\,(H P_\kappa)^{-1},
\qquad
(H P_\kappa)^2=(-1)^q I.
\]
因此 \(\langle \mathsf C_q,HP_\kappa\rangle\subset \mathrm{GL}(2^q,\ZZ)\) 为一类对 \(\langle \mathsf C_q\rangle\) 施加反演自同构的二面体型扩张；其在模 \(2\) 约化下退化为定理 \ref{thm:xi-boolean-binary-jump-kernels-char2-s3} 的 \(S_3\) 结构。
\end{proposition}
\begin{proof}
由定理 \ref{thm:xi-boolean-binary-jump-kernels-tensor-spectrum} 的逆核公式，
\[
\mathsf D_q^{-1}(u,v)=(-1)^{|u\wedge v|}\mathbf{1}_{\{u\vee v=\mathbf{1}\}},
\qquad
\mathsf C_q(u,v)=\mathbf{1}_{\{u\vee v=\mathbf{1}\}}.
\]
当 \(u\vee v=\mathbf 1\) 时有 \(|u|+|v|=|u\vee v|+|u\wedge v|=q+|u\wedge v|\)，故
\[
(-1)^{|u\wedge v|}
=(-1)^{|u|+|v|-q}
=(-1)^q\,(-1)^{|u|}\,(-1)^{|v|}.
\]
代回即得 \(\mathsf D_q^{-1}=(-1)^qH\mathsf C_qH\)。
同理由
\(
\mathsf C_q^{-1}(u,v)=(-1)^{\#\{i:\ u_i=v_i=0\}}\mathbf{1}_{\{u\wedge v=\mathbf{0}\}}
\)
与
\(
\mathsf D_q(u,v)=\mathbf{1}_{\{u\wedge v=\mathbf{0}\}}
\)
得到 \(\mathsf C_q^{-1}=(-1)^qH\mathsf D_qH\)。
\par
由互补共轭 \(\mathsf D_q=P_\kappa\mathsf C_qP_\kappa^{\mathsf T}\) 得
\(
\mathsf C_q^{-1}=(-1)^qH P_\kappa\,\mathsf C_q\,P_\kappa^{\mathsf T}H
=(-1)^q(HP_\kappa)\mathsf C_q(HP_\kappa)^{-1}
\)。
最后，因 \(|\kappa(u)|=q-|u|\)，有 \(P_\kappa H P_\kappa=(-1)^q H\)，从而
\[
(HP_\kappa)^2=H(P_\kappa H P_\kappa)=(-1)^qH^2=(-1)^qI.
\]
证毕。
\end{proof}

\begin{proposition}[\(2\)-进半周期标量化与张量幂的同余坍塌]\label{prop:xi-boolean-binary-jump-kernels-2adic-halfperiod}
\leanverified{paper\_xi\_boolean\_binary\_jump\_kernels\_2adic\_halfperiod}
沿用定理 \ref{thm:xi-boolean-binary-jump-kernels-tensor-spectrum} 的记号，并令
\(
A=\begin{psmallmatrix}0&1\\[2pt]1&1\end{psmallmatrix}
\)
从而 \(\mathsf C_q=A^{\otimes q}\)。
则对任意整数 \(k\ge 3\) 有：
\begin{enumerate}
\normalfont
  \item \textbf{基矩阵的半周期标量化：}
  \[
  A^{3\cdot 2^{k-2}}\ \equiv\ (1+2^{k-1})I_2\pmod{2^k}.
  \]
  \item \textbf{张量幂的标量化继承：}
  \[
  \mathsf C_q^{\,3\cdot 2^{k-2}}
  \ \equiv\
  (1+2^{k-1})^q\,I_{2^q}
  \ \equiv\
  \begin{cases}
  I_{2^q},& q\ \text{为偶数},\\
  (1+2^{k-1})I_{2^q},& q\ \text{为奇数},
  \end{cases}
  \pmod{2^k}.
  \]
\end{enumerate}
\end{proposition}
\begin{proof}
当 \(k=3\) 时，
\(
A^6=\begin{psmallmatrix}5&8\\[2pt]8&13\end{psmallmatrix}\equiv 5I_2=(1+2^2)I_2\pmod 8
\)，成立。
设对某 \(k\ge 3\) 已有
\(
A^{3\cdot 2^{k-2}}=(1+2^{k-1})I_2+2^kM
\)
（\(M\in\Mat_2(\ZZ)\)），则两边平方得
\[
A^{3\cdot 2^{k-1}}
=(1+2^{k-1})^2I_2+2^{k+1}(\cdots)
\equiv (1+2^k)I_2\pmod{2^{k+1}},
\]
其中用到 \(k\ge 3\) 时 \((1+2^{k-1})^2=1+2^k+2^{2k-2}\equiv 1+2^k\pmod{2^{k+1}}\)。
归纳完成 (1)。
\par
对 (2)，由 \(\mathsf C_q=A^{\otimes q}\) 得
\[
\mathsf C_q^{\,3\cdot 2^{k-2}}=\bigl(A^{3\cdot 2^{k-2}}\bigr)^{\otimes q}
\equiv \bigl((1+2^{k-1})I_2\bigr)^{\otimes q}=(1+2^{k-1})^q I_{2^q}\pmod{2^k}.
\]
当 \(k\ge 3\) 时，二项展开给出
\(
(1+2^{k-1})^q\equiv 1+q\,2^{k-1}\pmod{2^k}
\)；
因此 \(q\) 偶时该项同余于 \(1\)，而 \(q\) 奇时同余于 \(1+2^{k-1}\)。
合并即得结论。证毕。
\end{proof}

\begin{conjecture}[\(2\)-进阶的奇偶二分：\texorpdfstring{$\mathrm{ord}_{2^k}(\mathsf C_q)$}{ord} 的精确公式]\label{conj:xi-boolean-binary-jump-kernels-2adic-order-parity}
沿用命题 \ref{prop:xi-boolean-binary-jump-kernels-2adic-halfperiod} 的记号。
对 \(k\ge 3\) 令 \(\mathrm{ord}_{2^k}(\mathsf C_q)\) 表示 \(\mathsf C_q\) 在群 \(\mathrm{GL}(2^q,\ZZ/2^k\ZZ)\) 中的阶。
猜想存在仅由 \(q\) 的奇偶性决定的精确公式：
\[
\mathrm{ord}_{2^k}(\mathsf C_q)=
\begin{cases}
3\cdot 2^{k-1},& q\ \text{为奇数},\\[2pt]
3\cdot 2^{k-2},& q\ \text{为偶数}.
\end{cases}
\]
\end{conjecture}

\endinput


\begin{corollary}[奇异环 Cauchy 逆问题的唯一可逆性：测度密度的显式重构]\label{cor:xi-joukowsky-ellipse-cauchy-inversion-rho}
沿用定理 \ref{thm:xi-joukowsky-ellipse-plemelj-jump} 的设定，则 \(\rho\) 可由 \(G_{\sigma_L}\) 在椭圆上的跳跃显式恢复：
\[
\boxed{\;
\rho(\theta)
=-\sqrt{w(\theta)^2-4}\,
\Bigl(G_{\sigma_L}^{\mathrm{ext}}(w(\theta))-G_{\sigma_L}^{\mathrm{int}}(w(\theta))\Bigr)
\quad\text{a.e. }\theta.\;}
\]
因而映射 \(\rho\mapsto \sigma_L\)（等价地 \(\rho\mapsto G_{\sigma_L}\)）在 \(L^p,\ p>1\) 类上为单射；并且“给定椭圆外/内域解析函数并指定上述跳跃”给出一个标量 Riemann--Hilbert 问题，其解唯一。
\end{corollary}
\begin{proof}
由定理 \ref{thm:xi-joukowsky-ellipse-plemelj-jump} 直接代数消元即得重构公式。
唯一性由：若两解之差的跳跃为零，则差函数可解析延拓为整函数并在无穷远处衰减，从而恒为零。证毕。
\end{proof}

\begin{conjecture}[bin-fold 的 $\delta$-原子测度极限与平方根边缘跳跃]\label{conj:xi-foldbin-delta-measure-sqrt-edge-bridge}
对每个 \(m\ge 1\) 与 \(w\in X_m\)，令 \(\Delta_m(w)\) 为定理 \ref{thm:foldbin-uplift-unit-jump-delta-inversion} 中由 uplift 尾部偏移诱导的可达 \(\delta\)-集合。
取一族尺度 \(a_m\to\infty\) 及一族在 \(X_m\) 上的抽样机制（例如按 \(\pi_m(w)\propto d_m^{\mathrm{bin}}(w)\) 的推前分布，或限制到某一典型子族后取均匀），并定义归一化原子测度
\[
\mu_{m,w}:=\frac{1}{|\Delta_m(w)|}\sum_{\delta\in\Delta_m(w)}\delta_{\delta/a_m}.
\]
则存在某个绝对连续概率测度 \(\mu\) 使得在抽样意义下
\(\mu_{m,w}\Rightarrow \mu\) 弱收敛。
并且 \(\mu\) 的 Cauchy 变换 \(G_\mu\) 的边界跳跃在支撑端点呈现平方根型奇性：存在区间端点 \(A<B\) 与在 \((A,B)\) 上的密度 \(\varrho\) 使
\[
\dd\mu(x)=\varrho(x)\,\dd x,\qquad
\varrho(x)\asymp \sqrt{(x-A)(B-x)}\quad (x\to A^+,\,B^-),
\]
等价地，\(G_\mu\) 的 Plemelj 跳跃满足与定理 \ref{thm:xi-joukowsky-ellipse-plemelj-jump} 同型的平方根分母结构。
\end{conjecture}



\begin{theorem}[有限阿贝尔通道下的 Joukowsky 二值跳跃律完全解耦]\label{thm:xi-joukowsky-ellipse-jump-abelian-channel-decoupling}
令 \(H\) 为有限阿贝尔群，取 \(\mathcal{A}:=\CC[H]\cong\bigoplus_{\chi\in\widehat H}\CC\)。
设 \(\rho:[0,2\pi]\to\mathcal{A}\) 为可积密度，写 \(\dd\sigma(e^{\mathrm{i}\theta})=\rho(\theta)\frac{\dd\theta}{2\pi}\) 为 \(\TT\) 上的 \(\mathcal{A}\)-值绝对连续测度，
并按 Joukowsky 参数化 \(w(\theta)=r e^{\mathrm{i}\theta}+r^{-1}e^{-\mathrm{i}\theta}\)（\(L=r^2>1\)）推前得到 \(E_L\) 上的 \(\mathcal{A}\)-值测度 \(\sigma_L\)。
定义 \(\mathcal{A}\)-值 Cauchy 变换
\[
G_{\sigma_L}(w):=\int_{E_L}\frac{1}{w-\xi}\,\dd\sigma_L(\xi),\qquad w\in\CC\setminus E_L,
\]
并令 \(G^{\mathrm{ext}}_{\sigma_L},G^{\mathrm{int}}_{\sigma_L}\) 为椭圆两侧的非切向边界值（按分量逐点存在）。
则几乎处处 \(\theta\) 上成立与标量情形同形的跳跃公式
\[
\boxed{\;
G_{\sigma_L}^{\mathrm{ext}}\!\bigl(w(\theta)\bigr)-G_{\sigma_L}^{\mathrm{int}}\!\bigl(w(\theta)\bigr)
=-\frac{\rho(\theta)}{\sqrt{w(\theta)^2-4}}.\;}
\]
在 Fourier--Pontryagin 分解下，该恒等式对每个角色通道 \(\chi\in\widehat H\) 分量化为标量跳跃律
\[
\bigl(G_{\sigma_L}^{\mathrm{ext}}\bigr)_\chi\!\bigl(w(\theta)\bigr)-\bigl(G_{\sigma_L}^{\mathrm{int}}\bigr)_\chi\!\bigl(w(\theta)\bigr)
=-\frac{\rho_\chi(\theta)}{\sqrt{w(\theta)^2-4}},
\]
从而任何以该跳跃数据为约束的 \(\mathcal{A}\)-值二值跳跃问题（Cauchy/RHP）可完全等价地分解为 \(|\widehat H|\) 个彼此独立的标量二值跳跃问题。
\end{theorem}
\begin{proof}
由于 \(\mathcal{A}\cong\bigoplus_{\chi\in\widehat H}\CC\) 为有限维交换 \(\ast\)-代数，Cauchy 变换与非切向边界取值对取值空间线性，故
\[
G_{\sigma_L}=\bigoplus_{\chi\in\widehat H}\bigl(G_{\sigma_L}\bigr)_\chi,
\qquad
\rho=\bigoplus_{\chi\in\widehat H}\rho_\chi
\]
按分量逐一成立。
对每个 \(\chi\)，将定理 \ref{thm:xi-joukowsky-ellipse-plemelj-jump} 应用于标量密度 \(\rho_\chi\) 得到对应跳跃恒等式；再将这些标量恒等式组装回直和，即得 \(\mathcal{A}\)-值跳跃公式。
最后一段的“完全解耦”是上述分量化的直接语义重述。证毕。
\end{proof}

\begin{theorem}[椭圆调和测度下的对数平均闭式：奇异环 Mahler--Green 公式]\label{thm:xi-joukowsky-ellipse-mahler-green}
令 \(P(w)=a\prod_{j=1}^{n}(w-w_j)\in\CC[w]\) 为次数 \(n\) 的多项式，定义椭圆对数平均
\[
M_L(P):=\int_{E_L}\log|P(w)|\,\dd\nu_L(w)
=\frac{1}{2\pi}\int_{0}^{2\pi}\log\Bigl|P\!\left(r e^{\mathrm{i}\theta}+r^{-1}e^{-\mathrm{i}\theta}\right)\Bigr|\dd\theta,
\qquad r=\sqrt{L}.
\]
对每个根 \(w_j\) 任取分支并令
\[
u_{j,\pm}:=\frac{w_j\pm\sqrt{w_j^2-4}}{2},
\qquad
u_{j,+}u_{j,-}=1,
\qquad
z_{j,\pm}:=\frac{u_{j,\pm}}{r}.
\]
则有恒等式
\[
\boxed{\;
M_L(P)=\log|a|+\frac{n}{2}\log L+\sum_{j=1}^{n}\Bigl(\log\max\{1,|z_{j,+}|\}+\log\max\{1,|z_{j,-}|\}\Bigr).\;}
\]
若记 \(K_L\) 为 \(E_L\) 所围成的闭椭圆域，则上式中每个括号项均非负，且仅在 \(w_j\notin K_L\) 时为正；并且该项可理解为 \(\CC\setminus K_L\) 外域 Green 函数在 \(w_j\) 的显式值。
\end{theorem}
\begin{proof}
令 \(G(z):=z^{n}P(r z+r^{-1}z^{-1})\)。
展开可见 \(G\in\CC[z]\)，\(\deg G=2n\)，且首项系数为 \(a r^{n}=a L^{n/2}\)。
当 \(z=e^{\mathrm{i}\theta}\) 时 \(|z^{n}|=1\)，故
\[
M_L(P)=\frac{1}{2\pi}\int_{0}^{2\pi}\log|G(e^{\mathrm{i}\theta})|\,\dd\theta.
\]
对任意多项式 \(Q(z)=b\prod_{k=1}^{m}(z-\alpha_k)\) 的 Jensen/Mahler 恒等式给出
\[
\frac{1}{2\pi}\int_{0}^{2\pi}\log|Q(e^{\mathrm{i}\theta})|\,\dd\theta
=\log|b|+\sum_{k=1}^{m}\log\max\{1,|\alpha_k|\}.
\]
将其应用于 \(Q=G\)，并注意 \(G(\alpha)=0\iff P(r\alpha+r^{-1}\alpha^{-1})=0\)，于是对每个 \(w_j\) 的两根满足
\(
r\alpha=u_{j,\pm}
\)，即 \(\alpha=z_{j,\pm}\)。
代回即得所示闭式；非负性与 \(w_j\in K_L\iff |u_{j,\pm}|\le r\iff |z_{j,\pm}|\le 1\) 的等价性给出括号项的势论解释。证毕。
\end{proof}

\begin{theorem}[奇异环缺陷的非负性、椭圆域根判别与 \texorpdfstring{$L\downarrow 1$}{L↓1} 的区间谱极限]\label{thm:xi-joukowsky-ellipse-defect-nonnegativity-interval-limit}
在定理 \ref{thm:xi-joukowsky-ellipse-mahler-green} 的设定下，定义奇异环缺陷
\[
\Delta_L(P):=M_L(P)-\log|a|-\frac{n}{2}\log L.
\]
则成立：
\begin{enumerate}
  \item \(\Delta_L(P)\ge 0\)，并且 \(\Delta_L(P)=0\) 当且仅当 \(\{w_1,\dots,w_n\}\subseteq K_L\)。
  \item 若 \(P\in\RR[w]\)，则
  \[
  \lim_{L\downarrow 1}\Delta_L(P)=0
  \quad\Longleftrightarrow\quad
  \{w_1,\dots,w_n\}\subseteq[-2,2].
  \]
\end{enumerate}
\end{theorem}
\begin{proof}
(1) 由定理 \ref{thm:xi-joukowsky-ellipse-mahler-green}，
\(
\Delta_L(P)=\sum_{j,\pm}\log\max\{1,|z_{j,\pm}|\}
\)
为非负项之和，故 \(\Delta_L(P)\ge 0\)。
并且 \(\Delta_L(P)=0\) 当且仅当对每个 \(j\) 有 \(|z_{j,\pm}|\le 1\)，等价于 \(|u_{j,\pm}|\le r\)。
注意 \(u_{j,\pm}\) 为方程 \(u+u^{-1}=w_j\) 的两根，而映射 \(\varphi(u)=u+u^{-1}\) 将闭环域 \(\{u:r^{-1}\le |u|\le r\}\) 的像恰为 \(K_L\)（边界 \(|u|=r\) 的像为 \(E_L\)），从而 \(|u_{j,\pm}|\le r\iff w_j\in K_L\)。

(2) 当 \(L\downarrow 1\) 时 \(r\downarrow 1\)，环域 \(\{u:r^{-1}\le |u|\le r\}\) 收缩到 \(|u|=1\)，而 \(\varphi(|u|=1)=[-2,2]\)。
若存在实根 \(w_j\notin[-2,2]\)，则方程 \(u+u^{-1}=w_j\) 有一根满足 \(|u|>1\)，从而存在 \(\varepsilon>0\) 使得当 \(1<r<1+\varepsilon\) 时仍有 \(|u|>r\)，进而对应项 \(\log\max\{1,|u|/r\}\) 保持正下界，故 \(\Delta_L(P)\) 不可能趋于 \(0\)。
反之若全部实根皆在 \([-2,2]\)，则对每个 \(j\) 有 \(|u_{j,\pm}|=1\)，从而对任意 \(r>1\) 皆有 \(|u_{j,\pm}|\le r\)，由 (1) 得 \(\Delta_L(P)\equiv 0\)，特别地极限为 \(0\)。证毕。
\end{proof}
